\chapter{\texttt{core.common} --- shared Hamiltonian, mixing, and probabilities}
\label{ch:core-common}

This chapter covers the physics infrastructure shared by \emph{every} flavour
scenario (the 3-flavour SM and its BSM extensions) and by \emph{every}
propagation medium. It is the conceptual core of the whole project: this is
where the factorisation trick that makes matter propagation cheap to compute
is defined.

\section{\texttt{potential.py} --- kinetic and matter potentials}

\modulehead{core/common/potential.py}{dimensionless (matter + kinetic)
potentials, GPU-first.}

\begin{theorybox}
The propagation Hamiltonian in natural units splits into a kinetic part
(vacuum oscillation) and a coherent matter-interaction part (the MSW
effect),
\begin{equation}
  H = H_{\rm kin} + H_{\rm mat}.
\end{equation}
\textbf{Kinetic term.} For a mass eigenstate $i$ with energy $E \gg m_i$, the
phase accumulated relative to the lightest mode after travelling a distance
$L$ is $\Delta m^2_{i1} L / (2E)$ (in natural units, then divided by
$\hbar c$ to convert to laboratory units). Introducing the evolution scale
$L_{\rm scale}$ (by default the Earth radius $R_\oplus$) and the
dimensionless coordinate $x=L/L_{\rm scale}$, the dimensionless kinetic
eigenvalue entering the Hamiltonian is
\begin{equation}
  \keyresult{k_i = L_{\rm scale}\,\frac{\Delta m^2_i}{2E}\,\frac{1}{\hbar c}}.
\end{equation}
\textbf{Charged- versus neutral-current matter potentials.} Coherent
forward scattering off the background medium proceeds through two distinct
weak-interaction channels, each contributing an independent term to
$H_{\rm mat}$ and computed by its own dedicated function:
\begin{itemize}
  \item \textbf{Charged current (CC).} Only $\nu_e$ ($\bar\nu_e$) can
    exchange a $W^\pm$ with the ambient electrons (there are no positrons,
    muons, or taus to scatter off in ordinary matter), so this channel
    singles out the electron flavour alone. It adds an effective
    Wolfenstein-type potential \parencite{Wolfenstein1978},
    \begin{equation}
      V_{CC} = \sqrt{2}\, G_F\, n_e,
    \end{equation}
    which only affects the $ee$ entry of the flavour-basis Hamiltonian.
    Made dimensionless with the scale $L_{\rm scale}$ and computed by
    \pyfunc{matter\_potential\_cc},
    \begin{equation}
      \keyresult{V_{CC} = \pm\sqrt{2}\,G_F\,n_e\,L_{\rm scale}}, \qquad (+\text{ for } \nu,\ -\text{ for } \bar\nu).
    \end{equation}
  \item \textbf{Neutral current (NC).} A $Z^0$ is exchanged with electrons,
    protons, \emph{and} neutrons alike, and identically across all three
    \emph{active} flavours (lepton universality). In electrically neutral
    matter the proton and electron NC contributions cancel exactly, leaving
    \begin{equation}
      V_{NC} = -\frac{1}{\sqrt2}\,G_F\,n_n,
    \end{equation}
    where $n_n$ is the neutron density. Made dimensionless the same way and
    computed by \pyfunc{matter\_potential\_nc},
    \begin{equation}
      \keyresult{V_{NC} = \mp\frac{1}{\sqrt2}\,G_F\,n_n\,L_{\rm scale}}, \qquad (-\text{ for } \nu,\ +\text{ for } \bar\nu).
    \end{equation}
\end{itemize}
For antineutrinos both potentials flip sign, $V \to -V$ (the weak
interaction has opposite sign for particle/antiparticle), which is why both
functions take the same \code{antinu} flag. In the pure 3-flavour
Standard Model, $V_{NC}\,\mathbb{1}_3$ is a common phase across all three
active flavours and therefore has \emph{no effect on any oscillation
probability}: this is why $V_{NC}$ never appears at all in the SM matter
Hamiltonian, and why $n_n$ is nowhere plumbed through the SM code path. The
3+1 sterile extension changes this qualitatively, since the sterile flavour
is a gauge singlet and \emph{does not} receive $V_{NC}$: subtracting the
same global phase $V_{NC}\,\mathbb{1}_4$ from
$\diag(V_{CC}+V_{NC},\,V_{NC},\,V_{NC},\,0)$ leaves a genuine relative term
on the sterile diagonal entry,
\begin{equation}
  \diag(V_{CC}+V_{NC},\,V_{NC},\,V_{NC},\,0) - V_{NC}\,\mathbb{1}_4
  = \diag(V_{CC},\,0,\,0,\,-V_{NC}),
\end{equation}
which \emph{does} affect sterile-active oscillation probabilities, with
magnitude comparable to $V_{CC}$ itself: for isoscalar matter (equal proton
and neutron density, e.g.\ Earth's mantle) $V_{NC}/V_{CC}=n_n/(2n_e)
\approx1/2$, independent of the absolute density. See
\texttt{hamiltonian.py} below for how this term enters the reduced
Hamiltonian, and \cref{ch:earth,ch:atmosphere} for how $n_n$ is exposed from
the Earth and atmosphere density profiles.
\end{theorybox}

\begin{implbox}
\begin{itemize}
  \item \pyfunc{kinetic\_potential(DeltamSq, E, evolution\_scale\_m, ...)}
    implements $k = L_{\rm scale}\,\Delta m^2/(2E\hbar c)$, accepting scalar
    or tensor inputs with automatic \code{device}/\code{dtype} inference.
  \item \pyfunc{matter\_potential\_cc(n\_e, antinu, evolution\_scale\_m, ...)}
    implements $V_{CC}=\pm\sqrt2 G_F n_e L_{\rm scale}$; the sign is set by
    the boolean flag \code{antinu}.
  \item \pyfunc{matter\_potential\_nc(n\_n, antinu, evolution\_scale\_m, ...)}
    implements $V_{NC}=\mp\tfrac{1}{\sqrt2} G_F n_n L_{\rm scale}$ (same
    \code{antinu} sign convention as the CC potential). It has no
    \code{legacy\_precision} counterpart (see below), since the legacy
    \texttt{peanuts} reference never modelled the sterile extension.
  \item Public API input units: $n_e$, $n_n$ in $\mathrm{mol\,cm^{-3}}$,
    $\Delta m^2$ in $\mathrm{eV^2}$, $E$ in $\mathrm{MeV}$,
    \code{evolution\_scale\_m} in metres (default $R_\oplus$, see
    \code{util.constant.R\_E}).
\end{itemize}
\end{implbox}

\begin{notebox}
\textbf{The \code{legacy\_precision} flag.} Both \pyfunc{matter\_potential\_cc}
and \pyfunc{kinetic\_potential} accept a boolean \code{legacy\_precision}
keyword, but only \pyfunc{matter\_potential\_cc} does anything with it: it
selects which numerical value of the matter-potential prefactor
$\sqrt2\,G_F\,n_e$ is used.
\begin{itemize}
  \item \code{legacy\_precision=False} (the default): the prefactor
    $\sqrt2\,G_F\,N_A\times10^{6}\,(\hbar c)^2$ is computed once, at import
    time, from the full double-precision constants in
    \code{util/constant.py} (see the table below).
  \item \code{legacy\_precision=True}: the prefactor is replaced by the
    single hardcoded value $3.868\times10^{-7}\ \mathrm{MeV\,m^2}$,
    reproduced from Eq.~4.17 of the legacy \texttt{peanuts} reference
    implementation \parencite{LucentePeanuts}, itself only quoted to 4
    significant digits. This mode exists exclusively to reproduce
    \texttt{peanuts} output bit-for-bit in the legacy validation tests
    (\code{medium/*/test/*legacy*}); it is \textbf{not} intended for
    physics production use, since it discards precision that is otherwise
    available at no cost.
  \item \pyfunc{kinetic\_potential} accepts the same keyword purely for
    call-site symmetry with \pyfunc{matter\_potential} (so both can be
    driven by one shared \code{legacy\_precision} switch); no
    legacy-rounded kinetic prefactor exists, so the flag is a no-op there.
\end{itemize}

\begin{center}
\small
\begin{tabular}{@{}llll@{}}
\toprule
Symbol & Description & Value & Units \\
\midrule
$G_F$      & Fermi coupling constant     & $1.1663787\times10^{-11}$ & $\mathrm{MeV^{-2}}$ \\
$N_A$      & Avogadro constant           & $6.02214076\times10^{23}$ & $\mathrm{mol^{-1}}$ \\
$\hbar$    & Reduced Planck constant     & $6.582119569\times10^{-22}$ & $\mathrm{MeV\,s}$ \\
$c$        & Speed of light (exact)      & $2.99792458\times10^{8}$  & $\mathrm{m\,s^{-1}}$ \\
$\hbar c$  & derived from $\hbar$, $c$   & $1.9732698\times10^{-13}$ & $\mathrm{MeV\,m}$ \\
\midrule
\multicolumn{4}{@{}l}{Matter-potential prefactor $\sqrt2\,G_F\,N_A\times10^6(\hbar c)^2$:} \\
\midrule
full precision & derived from the constants above & $3.86793\times10^{-7}$ & $\mathrm{MeV\,m^2}$ \\
legacy (Eq.~4.17,~\parencite{LucentePeanuts}) & hardcoded, 4 sig.\ digits & $3.868\times10^{-7}$ & $\mathrm{MeV\,m^2}$ \\
\bottomrule
\end{tabular}
\end{center}

\textbf{Precision difference.} The two prefactors differ by a relative
$1.85\times10^{-5}$ ($\sim 15$ parts per million) --- five orders of
magnitude larger than double-precision machine epsilon
($\approx2.2\times10^{-16}$). This confirms the discrepancy is not a
floating-point rounding artefact but a deliberate consequence of the legacy
reference value carrying only 4 significant digits; \code{legacy\_precision}
therefore trades a small, well-quantified accuracy loss for exact
reproducibility of the legacy \texttt{peanuts} numbers.
\end{notebox}

\begin{notebox}
\begin{itemize}
  \item \textbf{Ultrarelativistic kinetic term.} \emph{Formulation:} the
    mass-eigenstate energy is linearised as $E_i\approx E+\Delta m_i^2/(2E)$
    to obtain $k_i=L_{\rm scale}\,\Delta m_i^2/(2E\hbar c)$
    (\pyfunc{kinetic\_potential}). \emph{Justification and validity:}
    requires $E\gg m_i$ for every propagating mass eigenstate; for the
    solar/reactor/atmospheric energies simulated here
    ($E\gtrsim\mathrm{MeV}$) and sub-eV neutrino masses this holds to better
    than one part in $10^{12}$, so the linearisation is exact for every
    scenario in this project.
  \item \textbf{Coherent forward elastic scattering.} \emph{Formulation:}
    only the forward ($\theta=0$) elastic-scattering amplitude contributes
    to $V_{CC}$/$V_{NC}$ --- the standard Wolfenstein index-of-refraction
    argument \parencite{Wolfenstein1978}; incoherent (large-angle or
    inelastic) scattering is neglected entirely.
    \emph{Justification and validity:} holds whenever the neutrino's mean
    free path for incoherent scattering is far longer than the propagation
    baseline, true throughout the Sun, Earth, and atmosphere at the
    MeV--GeV energies simulated here.
  \item \textbf{Electrically neutral, unpolarised matter for $V_{NC}$.}
    \emph{Formulation:} $n_p=n_e$ is assumed when deriving
    $V_{NC}=-\tfrac1{\sqrt2}G_F n_n$, so the proton and electron NC
    contributions cancel exactly and only the neutron density survives; the
    medium is further assumed unpolarised and at rest, so no macroscopic
    spin- or bulk-velocity-dependent term is added to $H_{\rm mat}$.
    \emph{Justification and validity:} exact for ordinary,
    non-magnetised astrophysical matter --- the Sun, Earth, and atmosphere
    all satisfy local charge neutrality to extremely high precision, and
    carry no coherent spin polarisation on the scales relevant here.
  \item \textbf{Local, instantaneous density.} \emph{Formulation:} $V_{CC}$
    and $V_{NC}$ are evaluated from the density $n_e(x)$, $n_n(x)$ at the
    neutrino's instantaneous position, treating the medium as locally
    static and homogeneous on the scale of the local oscillation length.
    \emph{Justification and validity:} requires the density profile to vary
    slowly compared to that oscillation wavelength; the segmented-evolutor
    machinery (\texttt{evolutor.py} below, and the numerical and
    perturbative chapters, \cref{ch:core-numerical,ch:core-perturbative})
    exists precisely to control the residual error from this approximation
    by refining the segmentation where the density changes fastest.
\end{itemize}
\end{notebox}

\section{\texttt{pmns.py} --- shared mixing infrastructure}

\modulehead{core/common/pmns.py}{abstract class \pyclass{PMNS} and container
\pyclass{PMNSParams}, common to every flavour scenario.}

\begin{theorybox}
\subsection*{Mixing angles and the CP-violating phase}
The active-flavour sector is parametrised by three mixing angles and one CP
phase, all real (angles in radians) \parencite{ParticleDataGroup}:
\begin{itemize}
  \item $\theta_{12}$ (\textbf{solar} angle) sets the $(e,\mu)$-plane
    rotation; the mixing angle measured by solar and by KamLAND reactor
    data, and the largest of the three.
  \item $\theta_{13}$ (\textbf{reactor} angle) sets the $(e,\tau)$-plane
    rotation; the smallest angle, measured by short-baseline reactor
    experiments (Daya Bay, RENO, Double Chooz). Its non-zero value is what
    makes CP violation observable at all in the lepton sector.
  \item $\theta_{23}$ (\textbf{atmospheric} angle) sets the $(\mu,\tau)$-plane
    rotation; measured by atmospheric and long-baseline accelerator
    experiments, close to (possibly exactly) maximal, $\theta_{23}\approx\pi/4$.
  \item $\delta$ (\textbf{Dirac CP phase}) is the only source of CP
    violation carried by this parametrisation (Majorana phases, relevant
    only if neutrinos are Majorana particles, do not affect oscillation
    probabilities and are absent from it). $\delta$ has no physical,
    rephasing-invariant effect unless all three angles are non-zero; it
    enters oscillation probabilities through the Jarlskog invariant
    (\pyfunc{PMNS.jarlskog\_invariant}, see below).
\end{itemize}
These four scalars, together with a \pyclass{RuntimeContext} fixing the
target \code{device}/\code{dtype}, are exactly the contents of
\pyclass{PMNSParams} --- nothing else: the mass-squared splittings do not
enter the mixing matrix, so they are kept out of this container and live in
\pyclass{MassSpectrum} instead (next section).

\subsection*{Elementary rotation and phase generators}
Every mixing matrix in this project, SM or BSM, is assembled from two
elementary building blocks that \pyclass{PMNS} implements once, sized
generically by \code{self.n\_flavours} so the same formulas serve the
$3\times3$ active-only case and any larger active+sterile case (BSM
chapter). Restricted to the base $3\times3$ active sector, with the fixed
flavour order $(e,\mu,\tau)=(0,1,2)$, they read explicitly:
\begin{align}
  R_{12}(\theta_{12}) &=
    \begin{pmatrix}
      \cos\theta_{12} & \sin\theta_{12} & 0 \\
      -\sin\theta_{12} & \cos\theta_{12} & 0 \\
      0 & 0 & 1
    \end{pmatrix}
  & \text{rotation in the $(e,\mu)$ plane;}
  \\
  R_{13}(\theta_{13}) &=
    \begin{pmatrix}
      \cos\theta_{13} & 0 & \sin\theta_{13} \\
      0 & 1 & 0 \\
      -\sin\theta_{13} & 0 & \cos\theta_{13}
    \end{pmatrix}
  & \text{rotation in the $(e,\tau)$ plane;}
  \\
  R_{23}(\theta_{23}) &=
    \begin{pmatrix}
      1 & 0 & 0 \\
      0 & \cos\theta_{23} & \sin\theta_{23} \\
      0 & -\sin\theta_{23} & \cos\theta_{23}
    \end{pmatrix}
  & \text{rotation in the $(\mu,\tau)$ plane;}
  \\
  \Delta(\delta) &= \diag(1,\,1,\,e^{i\delta})
  & \text{CP phase, carried at the $\tau$ index.}
\end{align}
$R_{12}$, $R_{13}$, $R_{23}$ are always \textbf{real}: \pyclass{PMNS} never
attaches a phase to them (they call the shared generator with
\code{phase=None}); only $\Delta$ carries $\delta$, and only at position
$(\tau,\tau)$. This is a modelling choice of the parametrisation, not a
computational shortcut: the same underlying generator (\pyfunc{PMNS.\_rot})
accepts an arbitrary rotation plane $(i,j)$ and an arbitrary phase, which is
exactly what lets the 3+1 sterile extension reuse it unchanged for the
active-sterile rotations $R_{14}$, $R_{24}$, $R_{34}$ (each carrying its own
CP phase) without modifying this base class at all (BSM chapter). How
$R_{12}$, $R_{13}$, $R_{23}$, $\Delta$ combine into a full mixing matrix ---
the product order, and which rotations sit inside versus outside the block
that commutes with the matter potential --- is scenario-specific and is
given in the SM and BSM chapters.

\subsection*{Flavour, mass, and reduced bases}
Three bases appear throughout this project:
\begin{itemize}
  \item the \textbf{flavour basis} $\alpha=e,\mu,\tau(,s)$: the physical
    basis in which neutrinos are produced and detected, and the basis in
    which the matter potential is always diagonal. In the plain 3-flavour
    Standard Model, or the 3+1 sterile extension without its optional
    neutral-current term, only the electron entry is populated,
    $H_{\rm mat}=\diag(V_{CC},0,\dots,0)$; once that neutral-current term is
    active, $H_{\rm mat}=\diag(V_{CC},0,0,-V_{NC})$ carries a second,
    non-zero entry on the sterile diagonal (\texttt{potential.py} section
    above, \texttt{hamiltonian.py} section below) --- diagonal in every
    case, but not necessarily supported on the electron entry alone;
  \item the \textbf{mass basis} $i=1,\dots,n$: the eigenbasis of the
    free-propagation (kinetic-only) part of the Hamiltonian, in which the
    propagation phases $e^{-ik_ix}$ are defined;
  \item an intermediate \textbf{reduced basis}, introduced purely for
    computational convenience (\texttt{hamiltonian.py} section below).
\end{itemize}
Every concrete \pyclass{PMNS} scenario provides a unitary $U$
(\pyfunc{pmns\_matrix}) relating the flavour and mass bases: an operator
$A_{\rm flavour}$ and its mass-basis counterpart $A_{\rm mass}$ are related
by $A_{\rm flavour} = U\,A_{\rm mass}\,U^\dagger$. In addition, it
factorises this $U$ as
\begin{equation}
  U = O\,U_{\rm red},
\end{equation}
for a second unitary $U_{\rm red}$ (\pyfunc{reduced}, the \emph{reduced
mixing matrix}) and a third, scenario-specific unitary $O$
(\pyfunc{outer\_block}, the \emph{outer block}). This factorisation defines
the reduced basis: it is the basis reached from the flavour basis by
applying $O$ \emph{alone},
\begin{equation}
  A_{\rm flavour} = O\,A_{\rm reduced}\,O^\dagger,
  \qquad
  A_{\rm reduced} = O^\dagger\,A_{\rm flavour}\,O,
\end{equation}
implemented generically, for any scenario's $O$, by \pyfunc{flavour\_basis}
and \pyfunc{reduced\_basis} (see the implementation box below). $U_{\rm
red}$ is precisely the remaining rotation that carries a reduced-basis
operator the rest of the way to the true mass basis: combining
$A_{\rm flavour}=U\,A_{\rm mass}\,U^\dagger=O\,A_{\rm reduced}\,O^\dagger$
with $U=O\,U_{\rm red}$ gives the reduced-to-mass-basis transformation and
its inverse,
\begin{equation}
  A_{\rm reduced} = U_{\rm red}\,A_{\rm mass}\,U_{\rm red}^\dagger,
  \qquad
  A_{\rm mass} = U_{\rm red}^\dagger\,A_{\rm reduced}\,U_{\rm red},
\end{equation}
so the reduced basis coincides with the mass basis exactly when $O$ is the
identity (equivalently $U_{\rm red}=U$).
That $U$ factorises this way, and that \pyfunc{flavour\_basis}/%
\pyfunc{reduced\_basis} implement exactly the $O$-dressing above, is the
core.common contract every concrete \pyclass{PMNS} subclass must satisfy;
\emph{which} concrete $U$, $O$, $U_{\rm red}$ a given scenario provides is
scenario-specific and is given in the SM and BSM chapters.
\end{theorybox}

\begin{implbox}
\begin{itemize}
  \item \pyclass{PMNSParams}: immutable container for
    $\theta_{12},\theta_{13},\theta_{23},\delta$ and the
    \pyclass{RuntimeContext}.
  \item \pyclass{PMNS} (abstract \pyclass{torch.nn.Module}): builds the
    shared generators \pyfunc{R12}, \pyfunc{R13}, \pyfunc{R23},
    \pyfunc{Delta} above, sized generically by \code{self.n\_flavours} (3
    for the SM, 4 for the 3+1 sterile case), and implements the part of the
    public interface that is \emph{independent of the number of flavours
    and of the product structure}: \pyfunc{refresh}, \pyfunc{flavour\_basis},
    \pyfunc{reduced\_basis}, \pyfunc{select\_antinu} (see box below),
    \pyfunc{conjugate}, \pyfunc{transpose}, \pyfunc{dagger}, the
    reduced-matrix variants
    \pyfunc{reduced\_conjugate}/\pyfunc{reduced\_transpose}/%
    \pyfunc{reduced\_dagger}, and the rephasing-invariant
    \pyfunc{jarlskog\_invariant}.
  \item Concrete subclasses only need to supply \pyfunc{pmns\_matrix},
    \pyfunc{reduced}, and \pyfunc{outer\_block}: how the four generators
    above combine into the full and reduced mixing matrices is genuinely
    different per scenario, so those three methods remain abstract here.
  \item \pyfunc{flavour\_basis} / \pyfunc{reduced\_basis} apply the generic
    dressing $O(\cdot)O^\dagger$ / its inverse $O^\dagger(\cdot)O$, for
    \emph{any} scenario-specific outer block $O$ (\pyfunc{outer\_block}) and
    any reduced-basis operator (not only the Hamiltonian); why this
    reproduces the full-flavour result exactly is shown in the
    \texttt{hamiltonian.py} section below.
\end{itemize}
\end{implbox}

\begin{implbox}
\textbf{\pyfunc{select\_antinu(U, antinu)}.} Applies the standard
antineutrino prescription to a full or reduced mixing matrix: returns $U$
unchanged for neutrinos and its complex conjugate $U^*$ for antineutrinos
(the sign flip of the matter potential, $V\to-V$, is handled separately by
\pyfunc{potential.matter\_potential}, not here). \code{antinu} may be a
plain Python \code{bool} (applies globally) or a boolean tensor broadcast
against the batch dimensions of $U$ via \code{torch.where}, so a single
batched call can mix neutrinos and antineutrinos within the same tensor ---
e.g.\ when propagating an atmospheric flux that contains both. Every mixing
matrix exposed by \pyclass{PMNS} (\pyfunc{pmns\_matrix}, \pyfunc{reduced},
\pyfunc{outer\_block}) routes through this single function, so the
neutrino/antineutrino convention is enforced in exactly one place.
\end{implbox}

\section{\texttt{mass\_spectrum.py} --- mass-squared splittings}

\modulehead{core/common/mass\_spectrum.py}{abstract class \pyclass{MassSpectrum}:
the two measurable mass-squared splittings shared by every oscillation
model.}

\begin{theorybox}
Oscillation probabilities depend on mass-squared \emph{differences}, never
on absolute masses, so the mass sector is parametrised by two independent,
directly measurable splittings \parencite{ParticleDataGroup}:
\begin{itemize}
  \item $\Delta m^2_{21}\equiv m_2^2-m_1^2$ (\textbf{solar} splitting): by
    convention always positive ($m_2>m_1$); measured by solar and KamLAND
    reactor data, $\Delta m^2_{21}\approx7.4\times10^{-5}\ \mathrm{eV^2}$.
    Sets the ``slow'' oscillation frequency.
  \item $\Delta m^2_{3\ell}$ (\textbf{atmospheric} splitting):
    $\Delta m^2_{31}\equiv m_3^2-m_1^2$ if $m_3$ is the heaviest state
    (\textbf{normal ordering}, $\ell=1$), or
    $\Delta m^2_{32}\equiv m_3^2-m_2^2$ if $m_3$ is the lightest
    (\textbf{inverted ordering}, $\ell=2$); measured by atmospheric and
    long-baseline accelerator data,
    $|\Delta m^2_{3\ell}|\approx2.5\times10^{-3}\ \mathrm{eV^2}$. Sets the
    ``fast'' oscillation frequency.
\end{itemize}
These are exactly the two fields stored by \pyclass{MassSpectrum}: no
absolute mass scale, and no third independent splitting (since
$\Delta m^2_{32}=\Delta m^2_{31}-\Delta m^2_{21}$ by construction, only two
of the three pairwise splittings are independent).
\end{theorybox}

\begin{notebox}[title={Normal and Inverted Ordering}]
Whether $m_3$ is the heaviest or the lightest mass eigenstate --- the
\textbf{neutrino mass ordering} --- is not yet known experimentally, and is
not fixed by any mixing angle: it is encoded entirely in the \emph{sign} of
$\Delta m^2_{3\ell}$, which \pyclass{MassSpectrum} takes as a single signed
input rather than as a separate ordering flag. Concretely,
\pyfunc{difference\_vector\_base()} builds the reduced (global-phase-removed)
three-component vector
\begin{equation}
  \bigl(\Delta m^2_1,\ \Delta m^2_2,\ \Delta m^2_3\bigr)
  = \begin{cases}
      \bigl(0,\ \Delta m^2_{21},\ \Delta m^2_{3\ell}\bigr) & \Delta m^2_{3\ell} > 0\ \text{(normal ordering, NO)},\\[2pt]
      \bigl(-\Delta m^2_{21},\ 0,\ \Delta m^2_{3\ell}\bigr) & \Delta m^2_{3\ell} < 0\ \text{(inverted ordering, IO)}.
    \end{cases}
\end{equation}
Both branches use the same stored $\Delta m^2_{21}>0$ and the same
$|\Delta m^2_{3\ell}|$: only the sign of the single scalar
$\Delta m^2_{3\ell}$ passed in switches the ordering, there is no separate
boolean or string argument anywhere in the API. This is a common source of
confusion when porting values from sources that specify ordering
explicitly: the sign must be set on $\Delta m^2_{3\ell}$ itself before it
reaches \pyclass{MassSpectrum}. The batched form evaluates both branches and
selects per-element with \code{torch.where}, so a single tensor of
$\Delta m^2_{3\ell}$ values with mixed signs is handled correctly in one
call.
\end{notebox}

\begin{implbox}
\begin{itemize}
  \item \pyclass{MassSpectrum} (abstract \pyclass{dataclass}): stores
    $\Delta m^2_{21}$, $\Delta m^2_{3\ell}$ as its only fields.
  \item \pyfunc{difference\_vector\_base()}: concrete, shared by every
    scenario; builds the 3-component vector above.
  \item \pyfunc{difference\_vector()}: abstract; sizes the base vector to
    the concrete scenario's flavour count --- returned unchanged for the
    plain 3-flavour case (\pyclass{MassSpectrum\_SM}, SM chapter), with a
    fourth entry $\Delta m^2_{41}$ appended for the 3+1 sterile case
    (\pyclass{MassSpectrum\_BSM}, BSM chapter). This split exists because
    sizing is scenario-specific while the active-sector vector itself is
    not.
  \item Consumed by \pyfunc{kinetic\_eigenvalue\_vector} (\texttt{potential.py})
    to build the dimensionless kinetic phases $k_i$, and bundled into
    \pyclass{OscillationParameters} as the \code{mass\_spectrum} field
    (next section).
\end{itemize}
\end{implbox}

\section{\texttt{oscillation.py} --- bundled oscillation state}

\modulehead{core/common/oscillation.py}{\pyclass{OscillationParameters}: an
immutable container grouping \code{pmns}, \code{mass\_spectrum},
\code{antinu}, and the optional \code{nsi} coupling.}

\begin{implbox}
This introduces no new physics: it replaces the previous pattern of passing
\code{(pmns, DeltamSq21, DeltamSq3l, antinu)} as independent arguments
through the entire propagation stack. \pyclass{OscillationParameters} is a
passive domain container: it does not build itself, and it never imports the
concrete SM/BSM/preset implementations it bundles. Building a fully
assembled \pyclass{OscillationParameters} from a named preset --- choosing
the concrete PMNS/mass-spectrum implementation and reading the parameter
values --- is done by
\code{PropagationConfig.oscillation\_parameters\_from\_preset} in
\code{tpeanuts/config/propagation.py}, from the preset registry in
\code{tpeanuts/config/presets.py}: that composition layer is intentionally
kept outside this module, which stays a plain container.
\end{implbox}

\begin{center}
\small
\begin{tabular}{@{}llll@{}}
\toprule
Attribute & Type & Default & Role \\
\midrule
\code{pmns} & \pyclass{PMNS} & required & \parbox[t]{4.4cm}{Built mixing object (\pyclass{PMNS\_SM} or \pyclass{PMNS\_sterile}).} \\[2pt]
\code{mass\_spectrum} & \pyclass{MassSpectrum} & required & \parbox[t]{4.4cm}{Built splitting object (\pyclass{MassSpectrum\_SM} or \pyclass{MassSpectrum\_BSM}).} \\[2pt]
\code{antinu} & \code{bool}/\code{Tensor} & \code{False} & \parbox[t]{4.4cm}{Antineutrino selector ($U\to U^*$, $V\to-V$).} \\[2pt]
\code{nsi} & \pyclass{NSIConfig} or \code{None} & \code{None} & \parbox[t]{4.4cm}{Optional NSI coupling; \code{None} means the plain SM matter potential.} \\[2pt]
\code{preset\_name} & \code{str} & \code{"custom"} & \parbox[t]{4.4cm}{Name of the preset used to build this object.} \\[2pt]
\code{ordering} & \code{str} & \code{""} & \parbox[t]{4.4cm}{Mass-ordering label carried by the preset (\code{"NO"}/\code{"IO"}).} \\[2pt]
\code{label} & \code{str} & \code{""} & \parbox[t]{4.4cm}{Short identifier string carried by the preset.} \\[2pt]
\code{description} & \code{str} & \code{""} & \parbox[t]{4.4cm}{Human-readable description/reference carried by the preset.} \\[2pt]
\midrule
\multicolumn{4}{@{}l}{Derived BSM-extension flags (read-only \code{@property}, not stored fields):} \\
\midrule
\code{BSM\_extension\_sterile} & \code{bool} & from \code{pmns} & \parbox[t]{4.4cm}{True iff \code{pmns.n\_flavours == 4} (a 3+1 sterile object).} \\[2pt]
\code{BSM\_extension\_NSI} & \code{bool} & from \code{nsi} & \parbox[t]{4.4cm}{True iff \code{nsi is not None}.} \\[2pt]
\code{BSM\_extension} & \code{bool} & \parbox[t]{1.6cm}{OR of the two above} & \parbox[t]{4.4cm}{True iff \emph{any} BSM extension is active; the routing flag checked by \code{core.BSM}, \code{core.numerical}, and \code{core.perturbative} to pick the SM-only fast path or the general BSM path.} \\[2pt]
\bottomrule
\end{tabular}
\end{center}

\begin{notebox}
The three \code{BSM\_extension*} flags are \emph{derived}, never stored:
they are recomputed on every access from \code{pmns}/\code{nsi} rather than
set at construction time, so they can never drift out of sync with the
objects they describe (there is no way to end up with
\code{BSM\_extension\_sterile=True} while \code{pmns} is actually a
3-flavour \pyclass{PMNS\_SM}). The design rationale for keeping
\code{mass\_spectrum} separate from \code{pmns} is that the PMNS mixing
matrix does not depend on $\Delta m^2_{21}/\Delta m^2_{3\ell}$ (those
parameters only enter through the kinetic term, previous sections), so
storing them on \pyclass{PMNSParams} as well would duplicate the same
values.
\end{notebox}

\begin{notebox}[title={Default Configuration}]
Unless a different preset is requested, the reference configuration used
throughout this project's examples and tests is the NuFIT 5.2 (2022)
Standard Model best fit with normal ordering, registered as
\code{"\_SM\_NUFIT52\_NO"} in
\code{tpeanuts.config.presets.OSCILLATION\_PRESETS} \parencite{Esteban2020NuFit,NuFit52}.
Built via
\code{PropagationConfig.oscillation\_parameters\_from\_preset("\_SM\_NUFIT52\_NO")},
it populates \pyclass{OscillationParameters} as follows:

\begin{center}
\small
\begin{tabular}{@{}lll@{}}
\toprule
Attribute & Default value & Note \\
\midrule
$\theta_{12}$ & $33.41^\circ$ & solar angle \\
$\theta_{13}$ & $8.58^\circ$ & reactor angle \\
$\theta_{23}$ & $49.0^\circ$ & atmospheric angle \\
$\delta$ & $197.0^\circ$ & Dirac CP phase \\
$\Delta m^2_{21}$ & $7.41\times10^{-5}\ \mathrm{eV^2}$ & solar splitting \\
$\Delta m^2_{3\ell}$ & $+2.511\times10^{-3}\ \mathrm{eV^2}$ & atmospheric splitting (positive) \\
\code{antinu} & \code{False} & neutrino, not antineutrino \\
\code{ordering} & \code{"NO"} & normal ordering, set by the positive sign above \\
\code{BSM\_extension\_sterile} & \code{False} & \code{pmns} is 3-flavour, not sterile \\
\code{BSM\_extension\_NSI} & \code{False} & \code{nsi = None} \\
\code{BSM\_extension} & \code{False} & no BSM extension active \\
\bottomrule
\end{tabular}
\end{center}
\end{notebox}

\section{\texttt{hamiltonian.py} --- the propagation Hamiltonian}

\modulehead{core/common/hamiltonian.py}{the single Hamiltonian-assembly
module in the project: correct and complete for the 3-flavour SM, the 3+1
sterile extension, Non-Standard Interactions, or any combination, with no
separate SM-only implementation.}

\begin{theorybox}
\subsection*{General construction of the flavour Hamiltonian}
In the flavour basis (\texttt{pmns.py} section above), the propagation
Hamiltonian is, in full generality, the mass-basis kinetic term rotated
into the flavour basis by the full mixing matrix $U$, plus whatever matter
term $H_{\rm mat}$ the active scenario contributes:
\begin{equation}
  H_{\rm flavour} = U\,\diag(k_1,\dots,k_n)\,U^\dagger + H_{\rm mat}.
\end{equation}
This decomposition is always correct, for \emph{any} unitary $U$ and
\emph{any} Hermitian $H_{\rm mat}$, with no assumption on how $U$
factorises or on what $H_{\rm mat}$ concretely is: it is simply the
Schr\"odinger equation written in the flavour basis. What \emph{does}
depend on the scenario is the concrete content of $H_{\rm mat}$
(\texttt{potential.py} section above): plain 3-flavour Standard Model, or
3+1 sterile \emph{without} its optional neutral-current term,
$H_{\rm mat}=\diag(V_{CC},0,\dots,0)$, populated on the electron entry
alone; with Non-Standard Interactions and/or the sterile neutral-current
term active, $H_{\rm mat}$ acquires extra diagonal and/or off-diagonal
entries beyond the electron one (full expression below, in ``The 3+1
sterile neutral-current term'' subsection). Taken at face value, evolving
$H_{\rm flavour}$ is expensive in every case: since $U$ and $H_{\rm mat}$ do
not commute in general, $H_{\rm flavour}$ must be re-diagonalised as a
genuinely complex $n\times n$ matrix at every point along the path where
the electron (and, when the NC term is active, neutron) density changes.

\subsection*{The reduction trick: factorising out the outer block}
Every concrete \pyclass{PMNS} scenario additionally factorises its mixing
matrix as $U = O\,U_{\rm red}$ (\texttt{pmns.py} section above), where the
scenario-specific \textbf{outer block} $O$ (\pyfunc{outer\_block}) is
constructed, in every scenario implemented so far, to leave the electron
flavour entry fixed, and therefore to commute exactly with the
\emph{plain charged-current} matter term:
\begin{equation}
  O\,\diag(V_{CC},0,\dots,0)\,O^\dagger = \diag(V_{CC},0,\dots,0).
\end{equation}
This commutation is a structural property of $O$ alone, and holds for
\emph{any} scenario's $O$ and $U_{\rm red}$ --- but, as stressed below, only
for this plain CC-only piece of $H_{\rm mat}$, not for $H_{\rm mat}$ in
general. Whenever $H_{\rm mat}$ reduces to this plain CC-only form (i.e.\
neither NSI nor the sterile neutral-current term is active), the general
expression above simplifies to
\begin{align}
  H_{\rm flavour} &= U \diag(k_1,\dots,k_n)\, U^\dagger + \diag(V_{CC},0,\dots,0) \nonumber\\
  &= O\,U_{\rm red}\diag(k)\,U_{\rm red}^\dagger\,O^\dagger + O\,\diag(V_{CC},0,\dots,0)\,O^\dagger \nonumber\\
  &= O \underbrace{\Bigl[U_{\rm red}\diag(k)\,U_{\rm red}^\dagger + \diag(V_{CC},0,\dots,0)\Bigr]}_{\displaystyle H_{\rm red}}\, O^\dagger
  \;=\; O\, H_{\rm red}\, O^\dagger.
  \label{eq:H-factorization}
\end{align}
\textbf{Computational consequence (the peanuts ``reduction trick'').} The
$n$-flavour evolution problem reduces to: (1) solving the evolution of the
smaller, cheaper reduced Hamiltonian $H_{\rm red}$, obtaining
$S_{\rm red}(x)$; (2) \emph{dressing} the result once with the constant
matrix $O$ (which depends on neither density nor $x$),
$S_{\rm flavour}(x) = O\, S_{\rm red}(x)\, O^\dagger$. This is exactly what
\pyfunc{pmns.flavour\_basis} does: the spectral calculations that
diagonalise the propagation Hamiltonian only ever need to act on
$H_{\rm red}$, never on $H_{\rm flavour}$. This shortcut is \emph{not}
available once $H_{\rm mat}$ carries an NSI or sterile neutral-current
contribution, since neither commutes with $O$ in general (see below, and
the dedicated implementation box on the explicit BSM $H_{\rm mat}$
construction).

\subsection*{Assembling the reduced Hamiltonian}
Given the reduced mass-squared vector $(\Delta m^2_1,\dots,\Delta m^2_n)$
from \pyclass{MassSpectrum} (\texttt{mass\_spectrum.py} section above) and
the reduced mixing matrix $U_{\rm red}$ of the active \pyclass{PMNS}
scenario, the reduced Hamiltonian is
\begin{equation}
  \keyresult{H_{\rm red} = U_{\rm red}\,\diag(k_1,\dots,k_n)\,U_{\rm red}^{\dagger} + \diag(V_{CC},0,\dots,0)}.
\end{equation}
The matter term populates only the electron-flavour entry, by construction
of $H_{\rm mat}$ (\texttt{potential.py} section above): this holds for any
flavour count $n$, so no special handling is needed as $n$ grows --- every
additional flavour beyond the first simply carries a zero matter coupling.
When Non-Standard Interactions are active, the matter term
generalises to $V_{CC}\,(\diag(1,0,\dots,0)+\varepsilon)$ for a Hermitian
coupling matrix $\varepsilon$, built in the flavour basis and rotated back
into the reduced basis with $O^\dagger(\cdot)O$.

\subsection*{The 3+1 sterile neutral-current term, and why it breaks the
reduction trick}
Optionally supplying a neutron density $n_n$ for a 4-flavour \pyclass{PMNS}
scenario (\texttt{potential.py} section above) adds the relative
neutral-current term derived there to the flavour-basis matter Hamiltonian,
\begin{equation}
  H_{\rm mat} = \diag(V_{CC},0,0,0) + V_{CC}\,\varepsilon - \diag(0,0,0,V_{NC}),
\end{equation}
(the NSI term present only when $\varepsilon$ is set). Crucially, the outer
block $O$ of the 3+1 sterile scenario mixes the $\{\mu,\tau,s\}$ block
(\cref{ch:core-bsm}) and therefore does \emph{not} commute with a term that
singles out the sterile diagonal entry:
\begin{equation}
  O\,\diag(0,0,0,V_{NC})\,O^\dagger \neq \diag(0,0,0,V_{NC}).
\end{equation}
This means the reduction-trick shortcut of the previous subsection ---
skipping the $O^\dagger(\cdot)O$ rotation entirely when $H_{\rm mat}$ is
already diagonal in the reduced basis --- is only valid when \emph{neither}
NSI \emph{nor} the neutral-current term is active. Whenever either is
present, \pyfunc{hamiltonian\_matter\_reduced} always takes the general
path: build $H_{\rm mat}$ in the flavour basis (as above) and rotate it,
$H_{\rm mat}^{\rm red} = O^\dagger\,H_{\rm mat}\,O$. Omitting $n_n$ (the
default) recovers the CC-only sterile matter term used throughout the rest
of this thesis, unchanged; the term has no effect at all for a 3-flavour
\pyclass{PMNS} (there $V_{NC}$ is the unobservable global phase discussed in
\texttt{potential.py}), so passing $n_n$ to a 3-flavour scenario is silently
ignored rather than raising.
\end{theorybox}

\begin{implbox}
\begin{itemize}
  \item \pyfunc{kinetic\_eigenvalue\_vector(oscillation, E, ...)}
    (\code{core/common/potential.py}) converts the mass spectrum's
    \pyfunc{difference\_vector()} (previous section) into the dimensionless
    kinetic phases $k_i$ via \pyfunc{kinetic\_potential}.
  \item \pyfunc{hamiltonian\_kinetic\_reduced(...)} implements
    $U_{\rm red}\diag(k)\,U_{\rm red}^\dagger$ for any flavour count, always
    using the Hermitian conjugate.
  \item \pyfunc{hamiltonian\_matter\_reduced(oscillation, n\_e\_mol\_cm3, *,
    n\_n\_mol\_cm3=None, ...)} implements $\diag(V_{CC},0,\dots,0)$ (no
    rotation) when neither NSI nor a genuine sterile NC term is requested,
    or the general rotated construction above otherwise. The
    \code{n\_n\_mol\_cm3} keyword is optional and only meaningful for a
    4-flavour \code{oscillation.pmns}.
  \item \pyfunc{\_matter\_direction\_reduced(oscillation, *, antinu,
    include\_nc, context)} (module-private): factors the reduced matter
    Hamiltonian into $H_{\rm mat}^{\rm red}=V_{CC}\,P_{CC} + V_{NC}\,P_{NC}$,
    returning the two density-independent ``direction'' matrices
    $P_{CC},P_{NC}$ on their own. Shared by
    \pyfunc{hamiltonian\_matter\_reduced} and by the perturbative
    first-order correction (\cref{ch:core-perturbative}), which needs the
    direction matrices in isolation to sandwich the spectral projectors.
  \item \pyfunc{hamiltonian\_reduced(...)} / \pyfunc{hamiltonian\_flavour(...)}
    sum both terms directly from $n_e$ (and optional $n_n$) and the
    splittings; \pyfunc{hamiltonian\_flavour} additionally applies the
    $O(\cdot)O^\dagger$ dressing via \pyfunc{pmns.flavour\_basis}.
\end{itemize}
\end{implbox}

\begin{implbox}[title={Explicit construction of \texorpdfstring{$H_{\rm mat}$}{H\_mat} under BSM extensions}]
\pyfunc{hamiltonian\_matter\_reduced} never rebuilds $H_{\rm mat}$ term by
term at every call; it factorises it into position-\emph{independent}
``direction'' matrices and position-\emph{dependent} scalar potentials, via
\pyfunc{\_matter\_direction\_reduced}:
\begin{equation}
  H_{\rm mat}^{\rm red} = V_{CC}\,P_{CC} \; [+\; V_{NC}\,P_{NC}].
\end{equation}
Concretely, in code order:
\begin{itemize}
  \item \textbf{Neither NSI nor the sterile NC term requested} (the common
    case). Fast path: $P_{CC}=\diag(1,0,\dots,0)$ directly in the reduced
    basis, with \emph{no} rotation, since this direction is exactly
    $O$-invariant (\labelcref{eq:H-factorization} above); $P_{NC}$ is not
    built at all.
  \item \textbf{NSI and/or the sterile NC term requested.} General path:
    build the flavour-basis direction matrices
    $D_{\rm flavour}=\diag(1,0,\dots,0)+\varepsilon$ (with $\varepsilon=0$
    if NSI is off) and, only if the NC term is requested,
    $D_{n,\rm flavour}=-\diag(0,\dots,0,1)$ (the sterile entry); rotate both
    into the reduced basis with the scenario's outer block $O$,
    \begin{equation}
      P_{CC} = O^\dagger\,D_{\rm flavour}\,O,
      \qquad
      P_{NC} = O^\dagger\,D_{n,\rm flavour}\,O.
    \end{equation}
  \item \pyfunc{hamiltonian\_matter\_reduced} then evaluates the
    position-dependent scalar potentials --- $V_{CC}$ always, via
    \pyfunc{matter\_potential\_cc}, and, only when \code{n\_n\_mol\_cm3} is
    supplied to a 4-flavour scenario, $V_{NC}$ via
    \pyfunc{matter\_potential\_nc} --- and assembles
    $H_{\rm mat}^{\rm red}=V_{CC}\,P_{CC}\,[+\,V_{NC}\,P_{NC}]$ above.
\end{itemize}
This split is exactly what lets \pyfunc{evolutor\_first\_order}
(\cref{ch:core-perturbative}) reuse $P_{CC}$, $P_{NC}$ directly, sandwiching
them against the spectral projectors without recomputing the rotation at
every energy or density sample.
\end{implbox}

\begin{notebox}
The reduced kinetic term is always built as $U_{\rm red}\diag(k)\,U_{\rm
red}^\dagger$ (the Hermitian conjugate, never the plain transpose): this is
the only form for which \pyfunc{pmns.flavour\_basis} applied to it
reproduces $U_{\rm full}\diag(k)\,U_{\rm full}^\dagger$ exactly, since
$U_{\rm full}=O\,U_{\rm red}$ per \labelcref{eq:H-factorization}. Whenever
$U_{\rm red}$ happens to be real for a given scenario, the Hermitian
conjugate and the transpose coincide there numerically; the distinction
only becomes observable once $U_{\rm red}$ is genuinely complex. This is,
in fact, exactly what happened in the original \texttt{peanuts}
implementation \parencite{LucentePeanuts}: since it only ever considered
the plain 3-flavour Standard Model, its reduced mixing matrix was always
real, so using the plain transpose there was correct, not merely
convenient -- the two conventions were indistinguishable in that restricted
setting. It is only once \texttt{tpeanuts} generalises to a genuinely
complex $U_{\rm red}$ (the 3+1 sterile extension with $\delta_{14}\neq0$,
\cref{ch:core-bsm}) that the plain transpose silently becomes wrong, which
is why the code always uses the Hermitian conjugate unconditionally rather
than branching on whether $U_{\rm red}$ happens to be real.
\end{notebox}

\section{\texttt{evolutor.py} (generic) --- operator composition}

\modulehead{core/common/evolutor.py}{applying an evolutor to a state and
composing ordered segment evolutors.}

The concrete construction of a single segment's evolutor $S(x_{k+1},x_k)$
--- closed-form, perturbative, or fully numerical --- is model-dependent
and is given in the corresponding SM and BSM chapters
(\cref{ch:core-sm,ch:core-bsm,ch:core-numerical,ch:core-perturbative}); this
module only composes and applies whatever evolutors those chapters produce.

\begin{theorybox}
\subsection*{Composing segment evolutors}
If the propagation path is split into $n$ consecutive legs
$x_0 \to x_1 \to \dots \to x_n$, the composition property of the
position-dependent Schr\"odinger equation gives
\begin{equation}
  \keyresult{S(x_n, x_0) = S(x_n,x_{n-1})\cdots S(x_2,x_1)\,S(x_1,x_0)},
\end{equation}
i.e.\ the total evolution operator is the \emph{ordered} product (most
recent leg on the left) of the per-leg operators. This identity is what
allows any continuous density profile to be treated as a sequence of
locally solvable segments, each with its own closed-form or numerically
integrated evolutor.

\subsection*{Applying the evolutor to a state}
Once the total evolutor $S(x_n,x_0)$ is known, it is applied to an initial
state as a plain matrix--vector product,
\begin{equation}
  \keyresult{\psi(x_n) = S(x_n,x_0)\,\psi(x_0)},
\end{equation}
implemented directly by \pyfunc{apply\_evolutor\_to\_state}. What ``the
state'' $\psi$ represents, and how a transition probability is read off
from the result, splits into two cases:
\begin{itemize}
  \item \textbf{Coherent propagation.} $\psi(x_0)$ is a genuine flavour
    amplitude vector (a quantum superposition carrying relative phase,
    e.g.\ a single produced flavour, $\psi(x_0)=e_\alpha$). The transition
    amplitude is obtained by applying the formula above directly, and the
    transition probability to flavour $\beta$ follows immediately as
    \begin{equation}
      P_\beta = |\psi_\beta(x_n)|^2 = \bigl|(S\,\psi(x_0))_\beta\bigr|^2.
    \end{equation}
  \item \textbf{Incoherent propagation.} The initial ensemble is a
    classical mixture of \emph{mass} eigenstates with weights
    $w_i\geq0$, $\sum_i w_i=1$, carrying no relative phase between them
    (e.g.\ neutrinos that lost coherence over a very long baseline). There
    is no single amplitude vector $\psi(x_0)$ to apply $S$ to; instead each
    mass eigenstate is propagated and projected onto the final flavour
    separately, and the resulting \emph{probabilities} --- not amplitudes
    --- are summed with the classical weights,
    \begin{equation}
      P_\beta = \sum_i \bigl|(S\,U_{\rm PMNS})_{\beta i}\bigr|^2\, w_i.
    \end{equation}
\end{itemize}
Both formulas, together with the unitarity check that verifies them, are
derived and implemented in full in the \texttt{probability.py} section
below.
\end{theorybox}

\begin{implbox}
\begin{itemize}
  \item \pyfunc{apply\_evolutor\_to\_state(S, psi)}: applies a scalar or
    batched evolutor to a state vector, with automatic broadcasting.
  \item \pyfunc{compose\_segment\_evolutors(segments)}: performs the ordered
    product above; internally uses a binary-tree reduction
    (\code{util.math.tree\_reduce\_matmul}) to minimise the depth of the
    matrix-multiplication chain on GPU.
\end{itemize}
\end{implbox}

\section{\texttt{probability.py} --- from evolutors to observable probabilities}

\modulehead{core/common/probability.py}{the single transition-probability
convention for the whole project.}

\begin{theorybox}
Given the flavour evolutor $S$ (row = final flavour, column = initial
flavour),
\begin{equation}
  P(\nu_\alpha\to\nu_\beta) = \bigl|S_{\beta\alpha}\bigr|^2.
\end{equation}
For a coherent initial state $\psi_\alpha$ (flavour superposition with
physical phase), the final flavour probability is
\begin{equation}
  P_\beta = \Bigl|\sum_\alpha S_{\beta\alpha}\psi_\alpha\Bigr|^2
          = \bigl|(S\psi)_\beta\bigr|^2.
\end{equation}
For an incoherent mass mixture with weights $w_i\geq0$, $\sum_i w_i=1$
(no relative-phase information between mass eigenstates, the regime selected
by \code{massbasis=True} throughout the code),
\begin{equation}
  P_\beta = \sum_i \bigl|(S\,U_{\rm PMNS})_{\beta i}\bigr|^2\, w_i .
\end{equation}
Probability conservation: since $S$ is unitary, $\sum_\beta P_{\beta\alpha}=1$
for every column $\alpha$; this is used as a numerical consistency check,
not as an additional physics ingredient.
\end{theorybox}

\begin{implbox}
\begin{itemize}
  \item \pyfunc{probability\_transition(S, ...)}: the full probability
    matrix, or a selected channel (transition or survival).
  \item \pyfunc{probability\_coherent\_state(psi)} / \pyfunc{probability\_coherent(S, psi)}:
    implement the coherent formula above.
  \item \pyfunc{probability\_incoherent(...)}: implements the incoherent
    formula, either from a precomputed probability matrix or by projecting
    mass weights through $S$ and $U_{\rm PMNS}$.
  \item \pyfunc{probability\_state(...)}: dispatches between both
    modes depending on the type of initial state received. Every medium's
    \code{[medium]\_probability\_state} (\cref{ch:vacuum,ch:solar,ch:earth,ch:atmosphere})
    delegates to this function.
  \item \pyfunc{normalize\_probability\_columns(P, eps=1e-15)}: rescales
    each column of a probability matrix to sum exactly to one, for callers
    that need a strictly conserved matrix after a numerically imperfect
    construction (e.g.\ a truncated spectral sum); \code{eps} guards the
    division against a vanishing column sum.
  \item \pyfunc{check\_probability\_conservation} / \pyfunc{check\_probability\_matrix}:
    numerical unitarity checks, used extensively in the tests
    (\code{core/BSM/test}, \code{core/common/test}).
\end{itemize}
\end{implbox}

\begin{theorybox}
\subsection*{Averaging a probability over a coordinate}
A probability is dimensionless and normalised, so averaging it over any
coordinate (energy, an arbitrary production variable, or the full zenith
range) must \emph{preserve} that normalisation: the natural reduction is a
\textbf{weighted average}, not a plain integral. Given a weight $w$ over a
grid $x$ (an explicit production spectrum for energy, an arbitrary
caller-supplied weight for any other coordinate, or the geometric
solid-angle element for angle),
\begin{equation}
  \keyresult{\langle P\rangle = \frac{\int P(x)\,w(x)\,\mathrm{d}x}{\int w(x)\,\mathrm{d}x}},
\end{equation}
evaluated by the trapezoidal rule, so $\langle P\rangle$ stays in $[0,1]$
regardless of how $w$ is normalised. This is the defining difference from
\texttt{flux.py} below: a flux is not itself normalised, so its energy
reduction is a plain (unnormalised) integral, not a weighted average. Like
every other function in this module, all three reductions below accept a
final flavour dimension of 3 or 4, the latter for the 3+1 sterile
extension.
\end{theorybox}

\begin{implbox}
Three functions implement the weighted-average pattern above, one generic
and two coordinate-specific:
\begin{center}
\small
\begin{tabular}{@{}lll@{}}
\toprule
Function & Weight $w$ & Formula implemented \\
\midrule
\parbox[t]{5.0cm}{\pyfunc{probability\_\allowbreak weighted\_\allowbreak average(\allowbreak P, coordinate, weight, dim=-2)}} &
  \parbox[t]{2.6cm}{arbitrary, caller-supplied} &
  \parbox[t]{5.6cm}{$\displaystyle \langle P\rangle=\frac{\int P(x)\,w(x)\,dx}{\int w(x)\,dx}$
  --- the generic primitive, used e.g.\ for production-height averages
  (\cref{ch:atmosphere})} \\[10pt]
\parbox[t]{5.0cm}{\pyfunc{probability\_integrated(\allowbreak P, E\_grid\_MeV, spectrum, energy\_dim=-2)}} &
  \parbox[t]{2.6cm}{production spectrum $w(E)$, \textbf{required} (no default)} &
  \parbox[t]{5.6cm}{$\displaystyle \langle P\rangle=\frac{\int P(E)\,w(E)\,dE}{\int w(E)\,dE}$} \\[10pt]
\parbox[t]{5.0cm}{\pyfunc{probability\_\allowbreak integrated\_\allowbreak angular(\allowbreak P, theta\_deg, angular\_dim=-2)}} &
  \parbox[t]{2.6cm}{$w(\theta)=\sin\theta$, fixed and geometric} &
  \parbox[t]{5.6cm}{$\displaystyle \langle P\rangle=\frac{\int P(\theta)\sin\theta\,d\theta}{\int \sin\theta\,d\theta}$} \\
\bottomrule
\end{tabular}
\end{center}
\pyfunc{probability\_integrated} has no default \code{spectrum}: a flat
weight over an arbitrary energy grid is a modelling choice, not a safe
default, so every caller must pass one explicitly (an explicit flat tensor
if that is genuinely intended). \pyfunc{probability\_integrated\_angular}
takes nothing in its place, since its weight is purely geometric. Both
energy- and angle-specific functions are self-contained implementations
(their own trapezoidal-rule code), not thin wrappers around
\pyfunc{probability\_weighted\_average}; the latter exists as the shared
pattern for coordinates --- such as production height --- that do not
warrant a dedicated function of their own.
\end{implbox}

\section{\texttt{flux.py} --- from probabilities to fluxes}

\modulehead{core/common/flux.py}{multiplying final flavour probabilities by
flux normalisations and optional spectra.}

\begin{theorybox}
\subsection*{From probabilities to a differential flux}
If $\Phi_\alpha^{\rm prod}$ is the flux (or its normalised spectral shape)
of the produced flavour $\alpha$, the detected flux in flavour $\beta$ is
\begin{equation}
  \Phi_\beta^{\rm det}(E) = \Phi_\alpha^{\rm norm}\cdot f(E)\cdot P(\nu_\alpha\to\nu_\beta; E),
\end{equation}
where $f(E)$ is an optional normalised spectrum. This module knows nothing
about the physical origin of $\Phi^{\rm prod}$ or the medium traversed: it
is purely the final algebraic step, shared by \texttt{medium.vacuum.flux},
\texttt{medium.solar.flux}, \texttt{medium.earth.flux}, and
\texttt{medium.atmosphere.flux}. When the full transition matrix
$P_{\beta\alpha}$ is available rather than an already-selected final-flavour
probability, the equivalent matrix form
$\Phi_\beta^{\rm det}=\sum_\alpha P_{\beta\alpha}\,\Phi_\alpha^{\rm prod}$
is provided directly by \pyfunc{flux\_transition} (see below), mirroring
\pyfunc{probability\_incoherent}'s plain-matrix branch
(\texttt{probability.py} above).

\subsection*{From differential flux to event rate}
$\Phi_\beta^{\rm det}(E)$ above is a \emph{differential} flux: a density in
energy (and, when angularly resolved, in zenith/nadir angle), not yet a
rate. Obtaining a neutrino \textbf{rate} --- the quantity that, multiplied
by an effective area and an exposure time, gives an expected event count
--- requires integrating the energy dependence out:
\begin{equation}
  \keyresult{\Phi_\beta^{\rm rate} = \int \Phi_\beta^{\rm det}(E)\;dE},
\end{equation}
evaluated numerically by the trapezoidal rule over a finite energy grid.
When $\Phi_\beta^{\rm det}$ additionally carries an angular dependence
(e.g.\ the atmospheric or Earth-crossing zenith angle), integrating only
over energy keeps that angular resolution, giving a rate resolved by both
angle and flavour. A further integration over the full zenith range is
provided by this module too (\pyfunc{flux\_integrated\_angular} below), as
the geometric counterpart to energy integration; unlike
\texttt{probability.py} above, both reductions here are \textbf{unnormalised}
totals -- a flux is a physical rate density, not a probability, so nothing
should be divided out. Every function in this module, like every function in
\texttt{probability.py}, accepts a final flavour dimension of 3 or 4 (the
3+1 sterile extension).
\end{theorybox}

\begin{implbox}
\begin{itemize}
  \item \pyfunc{flux\_state(probability, flux, spectrum=None)}:
    multiplies and preserves the final-flavour dimension; if
    \code{spectrum} is given, it is applied as an element-wise
    multiplicative factor before reduction. Every medium's
    \code{[medium]\_flux\_state} (\cref{ch:vacuum,ch:solar,ch:earth,ch:atmosphere})
    delegates to this function.
  \item \pyfunc{flux\_transition(probability\_matrix, initial\_flux)}:
    the matrix form
    $\Phi_\beta=\sum_\alpha P_{\beta\alpha}\Phi_\alpha$ described above, for
    callers that hold the full transition matrix rather than a
    final-flavour probability vector.
\end{itemize}
Three functions implement the unnormalised-integral pattern above, one
generic and two coordinate-specific:
\begin{center}
\small
\begin{tabular}{@{}lll@{}}
\toprule
Function & Coordinate & Formula implemented \\
\midrule
\parbox[t]{5.2cm}{\pyfunc{flux\_integrated\_coordinate(\allowbreak flux, coordinate, dim=-2)}} &
  \parbox[t]{2.2cm}{arbitrary} &
  \parbox[t]{5.8cm}{$\displaystyle \Phi_{\rm out}=\int \Phi(x)\,dx$ --- the generic,
  unnormalised primitive} \\[10pt]
\parbox[t]{5.2cm}{\pyfunc{flux\_integrated(\allowbreak flux, E\_grid\_MeV, energy\_dim=-2)}} &
  \parbox[t]{2.2cm}{energy $E$} &
  \parbox[t]{5.8cm}{$\displaystyle \Phi^{\rm rate}=\int \Phi(E)\,dE$ --- delegates directly to
  \pyfunc{flux\_integrated\_coordinate}} \\[10pt]
\parbox[t]{5.2cm}{\pyfunc{flux\_integrated\_angular(\allowbreak flux, theta\_deg, angular\_dim=-2)}} &
  \parbox[t]{2.2cm}{zenith angle $\theta$} &
  \parbox[t]{5.8cm}{$\displaystyle \keyresult{\Phi_{\rm total} = 2\pi\int \frac{\mathrm{d}\Phi}{\mathrm{d}\Omega}\sin\theta\,\mathrm{d}\theta}$
  --- its own implementation (extra $2\pi$ factor), not a wrapper} \\
\bottomrule
\end{tabular}
\end{center}
\pyfunc{flux\_integrated}'s default \code{energy\_dim=-2} matches a plain
\code{(..., E, 3)} flux; an angle-resolved \code{(..., E, angle, 3)} flux is
integrated with \code{energy\_dim=-3}, preserving the angle axis in the
result. \pyfunc{flux\_integrated\_angular} assumes azimuthal symmetry and
is evaluated by the trapezoidal rule with an explicit $2\pi$ azimuthal
factor -- unlike \pyfunc{probability\_integrated\_angular}
(\texttt{probability.py} above), none of the three reductions here
normalise: a flux is a physical rate density, not a probability, so
nothing should be divided out.
\end{implbox}

\section{\texttt{neutrino.py} --- flavour index convention}

\modulehead{core/common/neutrino.py}{\pyclass{FLAVOUR\_TO\_INDEX} dictionary,
\pyfunc{flavour\_index} validator, and \pyfunc{flavour\_state} builder.}

\begin{implbox}
Fixes, once and for all, the flavour index convention used in every tensor
across the project: $e\to0,\ \mu\to1,\ \tau\to2,\ s\to3$. The sterile index
has no charged-lepton counterpart, but the 3+1 extension
(\pyclass{PMNS\_sterile}, \cref{ch:core-bsm}) fixes its position at index 3,
so it is included under a small set of aliases (\code{"s"}, \code{"sterile"},
\code{"nus"}, \code{"nu\_s"}) for callers that want to select or build it by
name instead of hand-indexing a length-4 tensor. \pyfunc{flavour\_index}
accepts these text aliases and validates out-of-range integer indices,
preventing silent indexing errors when building initial states;
\pyfunc{flavour\_state} returns the corresponding unit coherent-state
vector, sized to \code{n\_flavours} (3 by default, or 4 to select the
sterile flavour).
\end{implbox}
